\section{Introduction}
\label{sec:introduction}

In large healthcare organizations such as Philips France, the management of
medical equipment relies on a complex portfolio of maintenance contracts,
manufacturer warranties and warranty extensions. This management is critical:
it directly conditions the continuity of clinical services, the quality of
patient care, and the operational costs borne by the organization. The
contractual portfolio evolves continuously, driven by installations, renewals
and expirations, which makes coverage mastery particularly difficult at the
scale of a large installed base.

Despite the growing availability of data from information systems, existing
tools remain largely descriptive and fragmented. Current dashboards offer
partial visibility into contractual coverage, but they neither anticipate
risks associated with contract expirations nor prioritize actions in a
proactive manner. They also fail to provide explanations accessible to
business users, limiting their ability to interpret indicators, justify
decisions, and act quickly in complex operational environments. The result
is a growing need for intelligent solutions capable of transforming data into
actionable decisions while delivering clear, contextualized justifications.

We argue that this problem is naturally expressed as sequential
decision-making under uncertainty and is therefore well suited to a
reinforcement learning (RL) treatment. Specifically, we formalize it as a
\emph{Markov decision process} (MDP) in which the state describes the set of
managed pieces of equipment and their current contractual posture
(maintenance, manufacturer warranty, extension), the actions correspond to
coverage decisions (subscription, renewal, adjustment, prioritization),
transitions model the portfolio dynamics and contract expirations, and the
reward balances coverage cost against the penalty associated with uncovered
equipment and late interventions. In this framework, an RL agent learns a
long-horizon coverage policy that anticipates expirations and prioritizes
the actions with the greatest value for the organization.

Operationalizing this formulation raises several challenges. The environment
is only partially observable and must be reconstructed from heterogeneous
data sources; the action space combines discrete decisions over a large
number of pieces of equipment; the cost of erroneous exploration in
production precludes direct online learning, which motivates the use of
simulators and offline evaluation protocols; finally, the learned policy
must be interpretable to be adopted by the business teams responsible for
contract management.

The contributions of this work are as follows:
\begin{itemize}
    \item We formalize contractual coverage optimization for medical equipment
          as a Markov decision process, integrating maintenance contracts,
          manufacturer warranties and warranty extensions.
    \item We develop a simulation environment reproducing coverage, expiration
          and intervention dynamics calibrated on operational characteristics
          of the installed base.
    \item We train and evaluate a Proximal Policy Optimization agent and
          compare it with two representative baselines of current practice
          (a random policy and a threshold-based heuristic).
    \item We empirically show that the learned policy reduces the number of
          uncovered equipment steps while remaining within reach of a strong
          heuristic reference, and we analyze the conditions under which RL
          is expected to provide a clearer advantage.
\end{itemize}

The remainder of the paper is organized as follows.
Section~\ref{sec:related_work} reviews related work on maintenance
optimization, service-contract design, and reinforcement learning applied to
healthcare operations. Section~\ref{sec:methodology} details the MDP
formulation and the learning algorithm. Section~\ref{sec:experiments}
describes the experimental protocol and reports the empirical results.
Section~\ref{sec:discussion} discusses interpretations, limitations and
perspectives, and Section~\ref{sec:conclusion} concludes.
